% !TEX root = ../main.tex

\chapter*{体例说明}
\addcontentsline{toc}{chapter}{体例说明}

\section*{一、层次与阅读顺序}

每经五层（与 V1《研究译注本（大正卷序）》相同）：

\begin{enumerate}[label=\textbf{\arabic*.}]
  \item \textbf{【正文·仿罗什风】}——全书主文。
  \item \textbf{【今译意】}——现代汉语，与主文逐段平行。
  \item \textbf{【底本·求那跋陀罗译】}——汉译《杂阿含》原文，作文献对照。
  \item \textbf{【巴利平行与参考译文】}——《相应部》等编号及巴利／英译摘要。
  \item \textbf{【校勘与说明】}——据平行校正及重建依据。
\end{enumerate}

经题用\textbf{简体}：\textbf{（经号）短题}（如「（一）无常」）；此处经号为\textbf{大正经号}。
题下另标「学术序第×经 · 大正第×经」。
◇ 或「重建经」：汉本略，据巴利或同型经补出。

\section*{二、信达雅（本书定义）}

\begin{description}[style=nextline,leftmargin=3em]
  \item[信] 对巴利平行所显\textbf{早期教理}负责；汉本仅定位与术语参照。
  \item[达] 主文与今译意段数、因果链一致。
  \item[雅] 仿罗什删冗、四字节奏；\textbf{非}繁转简，\textbf{非}伪托罗什手笔。
\end{description}

\section*{三、建议引用}

引用本书某一经时，\textbf{经号须与大正经号一致}（可兼标学术序），平行经编号须与该经【平行】栏所标主据一致（不可混用他经编号）。

\begin{quote}\small
《杂阿含经·仿罗什风新译》研究译注本（经序重排），大正第 1 经，据《相应部》22.12 校正，2026.
\end{quote}

\noindent 若主平行在《增支部》等，则改标该经编号，例如：大正第 275 经，据《增支部》8.9 校正。
无可靠巴利专经平行者，不写「据……校正」，可标「据汉本」或参【校勘与说明】。
引用 ◇ 重建经时，请加「据平行经重建」。

\clearpage
